\documentclass{article}
\usepackage[margin=1in]{geometry}

\author{alexander}
\date{\today}
\title{parametric equations and polar coordinates}
\begin{document}
\maketitle



recall the following before reading this document: $\frac{x}{r} = \cos(\theta)$, $\frac{y}{r} = \sin(\theta)$, $\frac{y}{x} = \tan(\theta)$ and $x^2 + y^2 = r^2$\\

in rectangular coordinates, a point is represented by $(x, y)$. function $\subseteq$ relation. every curve can be described by a relation between x and y, but not every curve can be described by a single function of x.\\

consider this equation for a circle $x^2 + y^2 = r^2$. when we solve for y or x we obtain two functions for the entire curve. it would be nice if we had one equation that we coule manipulate to produce individual points from for the entire curve... so we introduce a new independent variable upon which the components of the point depend.\\

consider two points $A = (1, 2), B = (5, 6)$, $\Delta x = 4, \Delta y = 4$. let $t=0$ describe no movement at all and $t=1$ to describe all of the movement. $x = x_0 + \Delta x \cdot t, y = y_0 + \Delta y \cdot t$\\

now let us begin with the function $y = x^2$. let us consider both possibilities for our new independent variable: $t = y$ and $t = x$. if $t = y$, then solving for $x$ gives two possible values for each positive value of $t$. this means that a single value of $t$ can correspond to two different points on the parabola, so it does not give us one continuous path through the curve. if instead we let $t = x$, then each value of $t$ corresponds to exactly one point on the parabola. as $t$ increases, we move from the left side of the parabola through the vertex and toward the right side. thus, although both choices are based on the same original function, the choice of parameter affects how the points of the curve are associated with values of the parameter...so different parameterizations can produce the same curve; however, the way they are parameterized changes how the curve is traversed.\\

in polar coordinates, a point is represented by $(r,\theta)$, where $r$ is the distance from the origin and $\theta$ is the angle measured from the positive x-axis. unlike rectangular coordinates, where a point is described by $(x,y)$, polar coordinates describe a point by its distance and direction from the origin. a polar equation can place a constraint on $r$, $\theta$, or a relationship between them, while the other quantity remains free to vary.\\

consider the simple polar equations $r = a$, $\theta = \theta_0$, and $r = \theta$. in $r = a$, the distance from the origin is fixed at $a$, while $\theta$ is free to vary. as the angle changes, the point moves around the origin while remaining the same distance from it, producing a circle centered at the origin. in $\theta = \theta_0$, the angle is fixed while $r$ is free to vary. as the distance from the origin changes, the point moves along a line through the origin at the fixed angle $\theta_0$. finally, in $r = \theta$, neither quantity is fixed: as the angle increases, the distance from the origin increases with it, producing an archimedean spiral. by combining $r$ and $\theta$ in different ways using familiar functions, we can produce many other types of curves, such as limaçons, rose curves, and lemniscates.\\

\begin{enumerate}
    \item $r = 2a\cos(\theta) + 2b\sin(\theta)$
    \item $r^2 = 2ar\cos(\theta) + 2br\sin(\theta)$
    \item $x^2+y^2 = 2ax+2by$
    \item $x^2-2ax+y^2-2by=0$
    \item $(x-a)^2+(y-b)^2=a^2+b^2$
    \item center=$(a,b)$, radius=$\sqrt{a^2+b^2}$
\end{enumerate}

\vspace{1em}

\begin{enumerate}
    \item A full circle has area $\pi r^2$ and angle $2\pi$, so a sector with angle $\theta$ occupies the fraction $\frac{\theta}{2\pi}$ of the circle.
    
    \item $A = \frac{\theta}{2\pi}\pi r^2$
    
    \item $A = \frac{1}{2}r^2\theta$
    
    \item For a small slice of a polar curve, the angle is $d\theta$, so $dA = \frac{1}{2}r^2d\theta$.
    
    \item If $r=f(\theta)$, then $dA = \frac{1}{2}f(\theta)^2d\theta$.
    
    \item Adding all the slices from $\alpha$ to $\beta$ gives $A = \frac{1}{2}\int_{\alpha}^{\beta}f(\theta)^2d\theta$.
\end{enumerate}






now onto derivatives...i though i understood this before but maybe not so....like as far as $\frac{dy}{dx} \Rightarrow \frac{\frac{dy}{dt}}{\frac{dx}{dt}}$ but i failed to consider this new $\frac{dy}{dx} = \frac{f^{'}(\theta_0)\sin(\theta_0) + f(\theta_0)\cos(\theta_0)}{f^{'}(\theta_0)\cos(\theta_0) - f(\theta_0)\sin(\theta_0)}$ whats this?

arc legth and surface area!

\end{document}
